\documentclass[a4paper,11pt]{article}
\usepackage[utf8]{inputenc}
\usepackage[T1]{fontenc}
\usepackage[ngerman]{babel}
\usepackage{amsmath,amssymb}
\usepackage[margin=2.3cm]{geometry}
\usepackage{enumitem}
\usepackage{fancyhdr}
\usepackage{array}
\usepackage[table]{xcolor}
\usepackage{tikz}
\definecolor{bgdark}{HTML}{1E1E1E}
\definecolor{fgtext}{HTML}{E6E6E6}
\definecolor{gridcol}{HTML}{4A4A4A}
\pagecolor{bgdark}
\arrayrulecolor{fgtext}
\AtBeginDocument{\color{fgtext}}
\renewcommand{\headrule}{\color{fgtext}\hrule height \headrulewidth}
\newcommand{\karo}[1]{\par\noindent
  \begin{tikzpicture}
    \draw[step=5mm, gridcol, very thin] (0,0) grid (\linewidth,#1);
  \end{tikzpicture}\par}

\newcommand{\diff}[1]{\hfill\normalfont\textbf{[#1]}}
\newcommand{\aufg}[2]{\subsection*{Aufgabe #1 \diff{#2}}}
\renewcommand{\arraystretch}{1.2}

\pagestyle{fancy}
\fancyhf{}
\lhead{RADM -- Graphen}
\rhead{Übungsblatt 04}
\cfoot{\thepage}

\begin{document}

\begin{center}
  {\LARGE\bfseries Übungsaufgaben Graphen}\\[2pt]
  {\large Relationen, Algebra, Diskrete Mathematik -- Teil 5}\\[10pt]
\end{center}

\noindent
\begin{tabular}{@{}p{0.6\textwidth}@{}p{0.4\textwidth}@{}}
  \textbf{Name:} \hrulefill & \textbf{Datum:} \hrulefill
\end{tabular}

\vspace{4pt}
\noindent\rule{\textwidth}{0.4pt}
\vspace{4pt}

\noindent\textit{Themen: Grundbegriffe, Knotengrad, Adjazenzmatrix/-liste, Graphtypen,
Pfade/Zyklen, Euler/Hamilton, bipartite Graphen, kürzester Pfad (Dijkstra, A*),
TSP (Brute Force, Nearest Neighbour, 2-Opt).}

\vspace{6pt}

% ---------------------------------------------------------------
\aufg{1}{leicht}
Ungerichteter Graph $G=(V,E)$ mit $V=\{1,2,3,4,5\}$ und
$E=\{\{1,2\},\{1,3\},\{2,3\},\{3,4\},\{4,5\}\}$.
\begin{enumerate}[label=(\alph*)]
  \item Bestimme den Knotengrad jedes Knotens.
  \item Bilde die Summe aller Grade und vergleiche mit $2\cdot|E|$.
  \item Enthält $G$ einen Zyklus? Wenn ja, gib einen an.
\end{enumerate}
\karo{5cm}

% ---------------------------------------------------------------
\aufg{2}{leicht}
Für den Graphen aus Aufgabe~1:
\begin{enumerate}[label=(\alph*)]
  \item Gib die Adjazenzmatrix $A$ an.
  \item Gib die Adjazenzliste an.
  \item Ist $A$ symmetrisch? Begründe.
\end{enumerate}
\karo{6cm}

% ---------------------------------------------------------------
\aufg{3}{leicht}
\begin{enumerate}[label=(\alph*)]
  \item Wie viele Kanten hat der vollständige Graph $K_n$? Berechne für $n=6$.
  \item Wie viele Kanten hat ein Baum mit $n$ Knoten?
  \item Welchen Knotengrad hat jeder Knoten in einem Kreis $C_n$?
\end{enumerate}
\karo{4.5cm}

% ---------------------------------------------------------------
\aufg{4}{mittel}
Ungerichteter Graph mit $V=\{0,1,2,3\}$ und Kanten
$\{0,1\},\{1,2\},\{2,3\},\{3,0\},\{0,2\}$.
\begin{enumerate}[label=(\alph*)]
  \item Bestimme die Knotengrade.
  \item Existiert ein \emph{Euler-Zyklus}? (Kriterium: zusammenhängend, alle Grade gerade.)
  \item Existiert ein \emph{Euler-Pfad}? Wenn ja, gib einen an.
  \item Gib einen \emph{Hamilton-Zyklus} an.
\end{enumerate}
\karo{6cm}

% ---------------------------------------------------------------
\aufg{5}{mittel}
Zeige, dass der Graph mit $V=\{1,2,3,4,5,6\}$ und Kanten
$\{1,4\},\{1,5\},\{2,4\},\{2,6\},\{3,5\},\{3,6\}$ \emph{bipartit} ist,
indem du eine geeignete Partition der Knotenmenge angibst.
\karo{5cm}

% ---------------------------------------------------------------
\aufg{6}{schwer}
Gewichteter, ungerichteter Graph mit Knoten $A,B,C,D,E,F$ und Kanten:
\[ A\text{-}B{:}2,\; A\text{-}C{:}4,\; B\text{-}C{:}1,\; B\text{-}D{:}7,\;
   C\text{-}E{:}3,\; D\text{-}E{:}2,\; D\text{-}F{:}1,\; E\text{-}F{:}5. \]
Führe den \textbf{Dijkstra}-Algorithmus vom Startknoten $A$ aus.
Gib die kürzesten Distanzen zu allen Knoten sowie den kürzesten Pfad von $A$ nach $F$ an.
\karo{7cm}

% ---------------------------------------------------------------
\aufg{7}{schwer}
Im Graphen aus Aufgabe~6 sei $F$ der Zielknoten. Gegeben die (zulässige) Heuristik
\[ h(A){=}8,\ h(B){=}7,\ h(C){=}6,\ h(D){=}1,\ h(E){=}3,\ h(F){=}0. \]
\begin{enumerate}[label=(\alph*)]
  \item Worin unterscheidet sich $A^{*}$ vom Dijkstra-Algorithmus?
  \item Führe $A^{*}$ von $A$ nach $F$ aus (Priorität $p=d+h$) und gib den kürzesten Pfad an.
\end{enumerate}
\karo{6.5cm}

% ---------------------------------------------------------------
\aufg{8}{mittel}
\textbf{TSP -- Brute Force.} Vier Städte $A,B,C,D$ mit symmetrischen Distanzen:
\[ AB{=}2,\ AC{=}9,\ AD{=}1,\ BC{=}3,\ BD{=}8,\ CD{=}9. \]
Bestimme durch Ausprobieren aller Touren ab $A$ (mit Rückkehr nach $A$) die
kürzeste Rundreise und ihre Länge.
\karo{6cm}

% ---------------------------------------------------------------
\aufg{9}{mittel}
\textbf{Nearest-Neighbour-Heuristik.} Wende auf die Distanzen aus Aufgabe~8 die
Nearest-Neighbour-Heuristik mit Start in $A$ an. Welche Tour und Länge ergeben sich?
Vergleiche mit dem Optimum.
\karo{5.5cm}

% ---------------------------------------------------------------
\aufg{10}{mittel}
\textbf{2-Opt.} Erkläre einen 2-Opt-Schritt und wende ihn auf die
Nearest-Neighbour-Tour aus Aufgabe~9 an, um sie zu verbessern.
Welche Tour und Länge erhältst du?
\karo{6cm}

\end{document}
